\section{관련 연구}
\label{sec:related}

DriveLM, DriveVLM과 Alpamayo-R1은 장면 이해와 행동ㆍ궤적 예측을 언어 추론에 연결한다 \cite{sima2024drivelm,tian2024drivevlm,wang2025alpamayo}. 최근 연구는 생성 설명이 장면과 행동에 충실한지, 센서 교란이나 객체 제거에도 추론--행동 관계가 유지되는지를 평가한다 \cite{mayumu2026faithfulness,nguyen2026vladrivebench,priyadershi2026lostinfog,panda2026cvaa}. 이들의 주된 단위는 한 장면의 추론ㆍ행동 정렬 또는 입력 개입에 대한 모델 반응이다.

본 연구는 모델에 같은 질문을 다시 하거나 차량 행동을 교정하지 않는다. 이미 저장된 CoC를 시간순으로 읽어 이전 사건에서 시작된 요구, 그 요구가 끝나는 조건과 행동 순서를 검사한다. 따라서 한 장면의 설명이 맞는지 평가하는 연구를 대체하지 않고, 여러 주석을 이어 읽었을 때 서로 모순되지 않는지 확인하는 보완적 품질 점검에 해당한다.

입력이나 가정을 일부 바꾸어 결과 변화를 보는 반사실 추론 연구는 행동 분류를 수정하거나 구조화된 설명 레이블을 학습에 넣어 계획을 교정한다. 위반 사례를 새로 만들어 모델을 개선하는 연구도 있다. 이와 달리 본 연구는 새 행동이나 학습 레이블을 생성하지 않으며, 검사 규칙을 통과하지 못한 항목을 곧바로 위험 자료로 바꾸지 않는다. 출력은 어떤 요구ㆍ처리 방식ㆍ증거 때문에 해당 판정이 나왔는지 함께 기록한 검토 후보이다.

\uppaal{}은 시간에 따라 상태가 바뀌는 규칙을 자동으로 검사하는 도구이며, Scenic과 VerifAI는 자율 시스템의 시험 상황 생성과 위반 탐색을 지원한다 \cite{behrmann2004uppaal,fremont2019scenic,dreossi2019verifai}. 본 논문은 새로운 환경이나 제어 방법을 찾는 데 \uppaal{}을 쓰지 않고, 기록된 사건을 순서대로 재생하며 요구의 생성ㆍ유지ㆍ종료를 확인하는 데 사용한다. 따라서 기여는 차량과 환경의 반복 상호작용 전체에 대한 안전 보증이 아니라, 저장된 연속 주석을 자동 실행하고 위반 이유를 추적할 수 있는 검사 규칙으로 만든 데 있다.

정리하면 기존 연구는 주로 한 장면의 설명이 행동과 맞는지, 입력을 바꾸면 행동도 달라지는지, 또는 실제 실행 시스템이 안전 조건을 지키는지를 묻는다. 본 연구는 같은 장면에 저장된 여러 CoC를 이어 읽으며 앞선 요구가 언제 시작되고 계속되며 끝나는지를 묻는다. 관련 연구는 이 질문의 차이를 설명하는 범위로 한정하고, 본 실험에서 직접 비교하지 않은 방법 간 우열은 논하지 않는다.

\begin{table*}[t]
\centering
\bitablecaption{핵심 관련 연구군과 본 연구의 주된 평가 단위 및 주장 범위.}{Primary evaluation units and claim scopes of the closest research families and this work.}
\label{tab:related-position}
\scriptsize
\begin{tabularx}{\textwidth}{L{0.19\textwidth}L{0.27\textwidth}L{0.23\textwidth}Y}
\toprule
Research area & Main item studied & Typical test or input & Relation to this work \\
\midrule
CoC explanation quality and action response \cite{mayumu2026faithfulness,nguyen2026vladrivebench,priyadershi2026lostinfog,panda2026cvaa} & Match between reasoning and action in one scene & Run the model again with changed input or a removed object & Shows whether an explanation matches the scene and action \\
Reasoning-based model improvement & Alternative action or reasoning label & Change planning or training & Changes model behavior; this work only checks stored annotations \\
Automatic autonomy-system testing \cite{behrmann2004uppaal,fremont2019scenic,dreossi2019verifai} & A driving situation or a stated system rule & Search environments, simulate, or check a rule & Provides checking tools; this work replays fixed annotation events \\
This work & Requirements carried across stored CoCs & Build rules, retain prior requirements, and compare fixed cases & Checks annotation consistency and data sources, not physical safety \\
\bottomrule
\end{tabularx}
\end{table*}

표~\ref{tab:related-position}은 어떤 방법이 더 우수한지를 비교하기보다 각 연구가 무엇을 근거로 어떤 질문에 답하는지를 구분한다. 장면별 설명 연구가 ``설명이 장면과 행동에 맞는가''를 묻는다면, 본 연구는 ``한 사건의 요구가 다음 사건까지 남을 때 두 주석이 함께 성립할 수 있는가''를 묻는다. 또한 실제 환경과 차량 제어를 함께 실행하여 시험하는 연구와 달리, 본 연구는 저장된 문장과 기록만 검사한다. 따라서 제어 결과가 다시 차량과 환경에 영향을 주는 전체 과정의 안전을 보장하지 않는다.
